% ============================================================
% chapters/02-problem.tex -- פרק 2: הבעיה, שאלת המחקר והיעדים
% Content only. All design comes from template.cls.
% New commands used here (see work/template-needs.md):
%   table / tabular with \toprule \midrule \bottomrule (booktabs)
%   tikzpicture (figure is in figures/ch2-chain.tex)
% ============================================================

\chapter{הבעיה, שאלת המחקר והיעדים}
\label{ch:problem}

פרק זה מגדיר את מרחב הבעיה וגוזר ממנו את הדרישות של המחקר. איור~\ref{fig:problem-chain} מראה את השרשרת: מהמצב המצוי, דרך הפער, אל שאלת המחקר ואל מפרט הדרישות שלפיו ייבחן המחקר בסופו.

\begin{figure}[htbp]
  \centering
  \input{figures/ch2-chain}
  \caption{מהמצב המצוי אל מפרט הדרישות}
  \label{fig:problem-chain}
\end{figure}

\section{המצב המצוי: ויכוח בלי מדידה}
\label{sec:problem-current}

הוויכוח על תפקידה של הנצרות בנפילת הקיסרות המערבית נמשך מאז 410, אבל הוא מתנהל כמעט כולו בכלים איכותניים. כפי שהוצג בפרק~\ref{ch:intro}, גיבון ראה בנצרות גורם מחליש \cite{gibbon_decline}, ואילו חוקרים בני זמננו מדגישים את הלחץ החיצוני \cite{heather2005fall}, את הקריסה הכלכלית \cite{wardperkins2005fall} או את האקלים והמגפות \cite{harper2017fate}. לכל אחת מהעמדות יש בסיס רציני במקורות, ובכל זאת אין ביניהן הכרעה. לפי הטיעון המוצע בדוח, הסיבה העיקרית לכך אינה מחסור בידע, אלא שלוש מגבלות שיטתיות:

\begin{enumerate}
  \item \textbf{מקרה אחד בכל פעם.} רוב המחקרים שנכללו בסקירה עוסקים במערב בלבד, או במזרח בלבד, או בפרנקים בלבד. כשבוחנים מקרה אחד, אי אפשר להפריד בין השפעת הדת לבין כל שאר הגורמים שפעלו באותו מקום ובאותו זמן.
  \item \textbf{דוגמאות נבחרות.} חוקרים נוטים להביא את הדוגמאות שתומכות בטענתם. בספרות שנכללה בסקירה לא זוהתה רשימה אחידה של אירועים ומדדים שכל העמדות נבחנות מולה.
  \item \textbf{הגדרות לא מדידות.} מושגים כמו ״החלשה״, ״יציבות״ או ״השפעת הכנסייה״ נשארים עמומים. בלי הגדרה מדידה, אי אפשר לבדוק אם טענה נכונה או שגויה.
\end{enumerate}

גם הניסיונות הכמותיים הקיימים אינם עונים על השאלה. מחקרים על מאגר \en{Seshat} בודקים חברות מכל העולם ולאורך אלפי שנים, ברזולוציה גסה, ואינם מתמקדים ביחסים שבין הכנסייה למדינה \cite{turchin2018quantitative,turchin2022disentangling}. מנגד, מחקרים שמדדו את התנצרות העילית הרומית במספרים \cite{barnes1995statistics,salzman2002making} עסקו במערב בלבד, ולא קשרו את הממצאים ליציבות המדינה.

\section{הפער: שלושה מקרים, אותה דת, שלוש תוצאות}
\label{sec:problem-gap}

התקופה שבין המאה הרביעית למאה התשיעית מציעה הזדמנות נדירה למחקר השוואתי. שלוש מדינות אימצו את אותה דת, ולמרות זאת הגיעו לתוצאות הפוכות (טבלה~\ref{tab:three-cases}). מבנה כזה, שבו גורם אחד משותף לכל המקרים והתוצאה שונה, מאפשר לשאול איזה גורם \textbf{אחר} מבדיל ביניהם.

\begin{table}[htbp]
  \centering
  \caption{שלושת המקרים במחקר}
  \label{tab:three-cases}
  \begin{tabular}{p{28mm}p{28mm}p{55mm}p{29mm}}
    \toprule
    \textbf{מדינה} & \textbf{תקופה במחקר} & \textbf{הקשר לנצרות} & \textbf{התוצאה} \\
    \midrule
    רומא המערבית & \LR{312--476} & דת חדשה במוסדות ותיקים & נפלה \\
    רומא המזרחית & \LR{312--610} & אותה נקודת פתיחה & שרדה עד \LR{1453} \\
    הממלכה הפרנקית & \LR{481--843} & נוצרית כמעט מהקמתה & התחזקה, התפצלה ב־843 \\
    \bottomrule
  \end{tabular}
\end{table}

הכינויים ״מערב״ ו״מזרח״ מציינים כאן מסלולים אזוריים לצורך ההשוואה, ולא שתי מדינות עצמאיות לאורך כל התקופה. לפני חלוקת 395, וגם לאחריה באירועים בעלי תחולה משותפת, אותם שליטים ומקורות עשויים לשרת את שני המסלולים. ישויות אלה לא יטופלו כשתי חזרות עצמאיות. אירועים משותפים יסומנו בנפרד וייבחנו גם בניתוח רגישות המשמיט אותם, לפי כללי האיגום שבפרק~\ref{ch:alternatives}.

\begin{keybox}{פער הידע}
בספרות שנכללה בסקירה לא זוהה מחקר שבחן, באותם מדדים ובשיטה אחת, מה מבדיל בין שלוש מדינות נוצריות שהגיעו לתוצאות שונות. בפרט, לא זוהתה מדידה של מידת ההשתלבות בין הכנסייה למוסדות המדינה כגורם נפרד מעצם קבלת הדת.
\end{keybox}

הספרות על הפרנקים מצביעה על כיוון אפשרי. הבישופים בגאליה ניסחו אידיאלים של שלטון נוצרי עוד בעת העתיקה המאוחרת, והמלכים הפרנקים אימצו אותם \cite{moore2011sacred}; מבנה הדיוקסיות של הבישוף רמיגיוס סיפק לממלכה של כלוֹדוויג בסיס טריטוריאלי \cite{poveda2019clovis}. כלומר, אצל הפרנקים הכנסייה שימשה תשתית של המדינה. ייתכן שזה ההבדל המכריע, וייתכן שלא – וזו בדיוק השאלה שהמחקר בא לבדוק.

\section{שאלת המחקר וההשערות}
\label{sec:problem-question}

הפרויקט נשען על שאלה והשערות הניתנות לבחינה: \textbf{האם התנצרות הייתה קשורה להחלשה או לחיזוק של יציבות מדינית בשלושת מקרי המחקר, בין המאה הרביעית למאה התשיעית, ואילו מאפיינים של השתלבות הכנסייה במוסדות השלטון מבחינים בין שני הכיוונים?} השערה \en{H1} גורסת שרמת ההתנצרות כשלעצמה לא תספק הסבר יציב להבדלי היציבות לאחר התחשבות בגורמים חלופיים. השערה \en{H2} מציעה קשר חיובי בין השתלבות מוסדית של הכנסייה לבין יציבות, לעומת יחסי תחרות וסכסוך. השערת האפס \en{H0} היא שאין קשר מדיד בין מדד ההשתלבות לבין מדד היציבות בתכנון הנבדק. אי־דחיית השערת האפס לא תוכיח היעדר קשר, והפרכת \en{H2} לא תוכיח שגיבון צדק. התוצאות יפורשו לצד אי־הוודאות, הכיסוי הדל וההסברים החלופיים; המחקר התצפיתי אינו מזהה לבדו השפעה סיבתית של הדת.

\section{הגדרת המשתנים}
\label{sec:problem-variables}

כדי שההשערות יהיו ניתנות לבדיקה, כל מושג מתורגם למדדים שאפשר לאסוף מהמקורות. כל מדד נמדד בכל אחת משלוש המדינות, במקטעי זמן של 25 שנה. טבלה~\ref{tab:variables} מציגה את המדדים המוצעים.

\begin{table}[htbp]
  \centering
  \caption{המשתנים במחקר והמדדים שלהם (הצעה ראשונית)}
  \label{tab:variables}
  \begin{tabular}{p{32mm}p{62mm}p{48mm}}
    \toprule
    \textbf{משתנה} & \textbf{מדד} & \textbf{מקור הנתונים} \\
    \midrule
    \multicolumn{3}{l}{\textit{משתנה תלוי: יציבות המדינה}} \\
    & מספר ניסיונות תפיסת שלטון ומרידות & כרוניקות ותולדות \\
    & משך שלטון ממוצע של השליט & רשימות שליטים \\
    & שינוי בשטח שבשליטת המדינה & אטלסים היסטוריים, \en{Seshat} \\
    & מספר שנים של מלחמת אזרחים & כרוניקות ותולדות \\
    \midrule
    \multicolumn{3}{l}{\textit{משתנה מסביר: השתלבות כנסייה–מדינה}} \\
    & בישופים בתפקידי ממשל ושיפוט & מכתבים, החלטות מועצות \\
    & חקיקה דתית מטעם השליט & קודקסים, קפיטולריות \\
    & מועצות כנסייה שכינס השליט & פרוטוקולים של מועצות \\
    & סכסוכים גלויים בין שליט לבישופים & כרוניקות, מכתבים \\
    \midrule
    \multicolumn{3}{l}{\textit{משתנה מסביר: רמת ההתנצרות}} \\
    & שיעור העילית הנוצרית & מחקרים קיימים \cite{barnes1995statistics,salzman2002making} \\
    \midrule
    \multicolumn{3}{l}{\textit{משתני פיקוח}} \\
    & לחץ צבאי חיצוני (פלישות, קרבות) & כרוניקות, \en{Seshat} \\
    & אינדיקטורים כלכליים (מטבעות, מסים) & ממצאים ארכאולוגיים \\
    & אקלים ומגפות & מחקרי אקלים \cite{harper2017fate} \\
    \bottomrule
  \end{tabular}
\end{table}

את מדדי ההשתלבות אי אפשר לאסוף מטבלאות קיימות. הם מצויים בתוך אלפי טקסטים: מכתבים, החלטות של מועצות כנסייה, חוקים וכרוניקות. לכן חלק מרכזי במחקר הוא ניתוח של טקסטים אלה בעזרת מודלי שפה, עם תיקוף בידי מומחה (פרק~\ref{ch:alternatives}).

\section{מפרט הדרישות}
\label{sec:problem-requirements}

מפרט הדרישות הוא החוזה שלפיו ייבחן המחקר. לכל דרישה יש מזהה, ערך יעד ושיטת אימות (טבלה~\ref{tab:requirements}). דרישה בלי שיטת בדיקה היא משאלה, ולכן אף דרישה כאן אינה נשארת בלי דרך לבדוק אותה.

\begin{table}[htbp]
  \centering
  \caption{מפרט הדרישות של המחקר}
  \label{tab:requirements}
  \small
  \begin{tabular}{p{29mm}p{32mm}p{44mm}p{37mm}}
    \toprule
    \textbf{מזהה} & \textbf{הדרישה} & \textbf{ערך יעד} & \textbf{אימות} \\
    \midrule
    \en{REQ-DATA-01} & כיסוי מאגר הנתונים & 34 מקטעים לפחות, 12 משתנים בכל אחד & בדיקת שלמות אוטומטית \\
    \en{REQ-DATA-02} & תיעוד מקור לכל ערך & 100\% מהערכים & בדיקת שדות מלאה וביקורת תוכן מדגמית \\
    \en{REQ-TXT-01} & גודל קורפוס הטקסטים & 2,000 קטעים מתויגים לפחות & ספירה במאגר \\
    \en{REQ-AI-01} & דיוק הסיווג של המודל & התאמה של 85\% לפחות לתיוג מומחה & מדגם של 300 קטעים \\
    \en{REQ-AI-02} & עקביות המודל & שינוי של עד 5\% בין שתי הרצות & הרצה חוזרת על אותו מדגם \\
    \en{REQ-STAT-01} & בדיקת ההשערות & מובהקות 0.05, דיווח על גודל האפקט & סקריפט ניתוח פתוח \\
    \en{REQ-REP-01} & יכולת שחזור & כל התוצאות משוחזרות מהקוד והנתונים & הרצה נקייה במחשב אחר \\
    \bottomrule
  \end{tabular}
\end{table}

מספר המקטעים בדרישה \en{REQ-DATA-01} נגזר מהתקופות שבטבלה~\ref{tab:three-cases}: כ־7 מקטעים של 25 שנה במערב, כ־12 במזרח וכ־15 אצל הפרנקים.

\subsection*{אילוצים}

\begin{itemize}
  \item \textbf{זמן:} 12 חודשי עבודה, צוות של שני חוקרים.
  \item \textbf{שפות המקורות:} לטינית ויוונית. התרגום והסיווג נעשים בעזרת מודלי שפה, ומחייבים תיקוף אנושי.
  \item \textbf{שרידות המקורות:} רוב הטקסטים ששרדו נכתבו בידי אנשי כנסייה. הטיה זו מטופלת בפרק~\ref{ch:risks}.
  \item \textbf{עלות:} ניתוח אלפי טקסטים במודלי שפה כרוך בעלות של טוקנים, שתוגבל בתקציב מוגדר מראש.
\end{itemize}
